% PROBLEM FORMULATION. Numbers only via \R macros.

\subsection{BLOCKS, BUDGETS AND VALUE}
\begin{definition}[Typed context block]\label{def:block}
A block $b=(\tau_b,\{x_b^v\}_v,D_b)$ has a type $\tau_b$ from a fixed set $\mathcal{T}$ (goal, instructions, tool schema,
skill document, fact, example, history turn, note), one or more textual variants $x_b^v$ (\emph{full} and optionally
\emph{short}), and a set $D_b$ of blocks it depends on. Its cost $c_m(b,v)$ is the number of tokens of its rendered variant
under model $m$'s tokenizer.
\end{definition}

An instance $q$ of a task family $f$ is a set of blocks $\mathcal{B}_q$ together with a deterministic success check. A
\emph{selection} $S$ assigns to each block either a variant or nothing; it is \emph{admissible} under budget $B$ if every
mandatory block (goal and instructions) is included, every included block's dependencies are included, at most one
variant per block is chosen, and $\sum_{b\in S} c_m(b,S(b))\le B$. A fixed, documented renderer turns $S$ into one prompt,
so success $Y_m(q,S)\in\{0,1\}$ is a function of the selection and the model's decoding.

\begin{definition}[Marginal context value profile]\label{def:mcv}
For model $m$ and family $f$, the MCV of type $\tau$ is
$\mathrm{MCV}(\tau)=\mathbb{E}_q\big[Y_m(q,\mathcal{B}_q)-Y_m(q,\mathcal{B}_q\setminus\tau)\big]$,
the expected loss in success when every block of type $\tau$ is removed from the full context. The pairwise interaction
of types $\tau_1,\tau_2$ is
$I(\tau_1,\tau_2)=\mathbb{E}_q\big[Y(\mathcal{B})-Y(\mathcal{B}\setminus\tau_1)-Y(\mathcal{B}\setminus\tau_2)+
Y(\mathcal{B}\setminus\{\tau_1,\tau_2\})\big]$,
positive when the two types are complements and negative when they are redundant or in conflict. The short-variant value
$\mathrm{MCV}_{\mathrm{short}}(\tau)$ is defined like $\mathrm{MCV}(\tau)$ with type $\tau$ rendered short instead of removed.
A profile is the table of these quantities with confidence intervals.
\end{definition}

\subsection{THE ASSEMBLY PROBLEM}
Given a profile, an instance and a budget $B$, the compiler must choose an admissible selection that maximizes predicted
value. Type-level values say how much a \emph{kind} of block is worth; instances contain many blocks of one type (dozens of
tool schemas or documents), so the compiler also needs instance-level block values $\hat v_b$, predicted from cheap
features. The resulting problem is a knapsack with dependency constraints, variant choice and pairwise type-presence
terms; it is NP-hard in general, and with negative interactions the objective is not submodular, so greedy guarantees for
coverage objectives~\cite{pacms2026,skillselect2026} do not apply.

\subsection{WHAT IS AND IS NOT GUARANTEED}
Given its inputs, the compiler returns an admissible selection that is optimal for the predicted objective whenever the
solver reports optimality, and it fails loudly if the mandatory blocks alone exceed the budget; it never silently
truncates. Nothing guarantees that predicted value equals realized success: profiles are estimated from finite ablations
on a development split, block values are predicted from features, and transfer to another model is an empirical
question. These are exactly the claims Section~\ref{sec:eval} tests.
